\section{Switching 坐标与周期 fiber}
\label{sec:gauge}

式 \eqref{eq:fixed-minimum} 用边符号表述优化，但边符号并非不变量。我们先转到局部圈符号，
同时保留有限 Hamilton 圈上无法消去的一个全局符号。所得 word 的周期性随后给出全文所用
的有限 Bloch 分解。

\subsection{符号圈平方图与 Hamilton 规范}

记 step-one 边 $\{i,i+1\}$ 的符号为 $a_i$，step-two 边 $\{i,i+2\}$ 的符号为
$b_i$。switching 函数 $\varepsilon$ 作用为
\[
 a_i\mapsto\varepsilon_i a_i\varepsilon_{i+1},\qquad
 b_i\mapsto\varepsilon_i b_i\varepsilon_{i+2}.
\]
相应邻接矩阵由 $D_\varepsilon=\diag(\varepsilon_i)$ 作正交相似变换。以 $i$ 为起点的
三角形符号与 step-one Hamilton 圈符号分别为
\begin{equation}
 \tau_i=a_ia_{i+1}b_i,
 \label{eq:triangle-flux}
\end{equation}
\begin{equation}
 \alpha=\prod_{i=0}^{n-1}a_i.
 \label{eq:ham-holonomy}
\end{equation}
二者均在 switching 下不变，并共同决定 switching 类。沿切开的 Hamilton 路径递归选取
$\varepsilon_{i+1}=\varepsilon_i a_i$ 后，局部 step-one 符号均变为正，step-two 符号
变为 $\tau_i$，绕行一周后的不匹配正是 $\alpha$。

在 lift 上，这一规范写成
\begin{equation}
 x_{i+n}=\alpha x_i,
 \label{eq:finite-holonomy-boundary}
\end{equation}
以及
\begin{equation}
 (A_\tau x)_i=x_{i-1}+x_{i+1}
   +\tau_{i-2}x_{i-2}+\tau_i x_{i+2}.
 \label{eq:hamilton-operator}
\end{equation}
因此“正 step-one 规范”是切开 lift 上的局部归一化；负 holonomy 位于边界条件，而不是
说有限圈上所有正边的乘积可以为负。

\subsection{对称性与本原周期}
\label{sec:symmetries}

令
\begin{equation}
 Q_i=\tau_i\tau_{i+1}.
 \label{eq:q-word}
\end{equation}
$Q$ 在 $\tau\mapsto-\tau$ 下不变。$p$-周期 $Q$ 存在 $p$-周期 lift 当且仅当
$\prod_iQ_i=1$；固定 $\tau_0$ 后由 $\tau_{i+1}=Q_i\tau_i$ 得到两个 lifts。

设 $\mathcal B_z^{(p)}$ 为满足 $x_{i+p}=zx_i$ 的序列空间，$H_\tau^{(p)}(z)$ 为
\eqref{eq:hamilton-operator} 在该空间上的限制，并定义
\begin{equation}
 R_p(\tau)=\max_{|z|=1}\rho(H_\tau^{(p)}(z))^2.
 \label{eq:bloch-edge}
\end{equation}

\begin{proposition}
\label{prop:invariances}
$R_p$ 满足：$R_p(-\tau)=R_p(\tau)$；循环平移在同一 $z$ 上保持 fiber 谱；若
$\tau_i^\vee=\tau_{-i-2}$，则
\[
 \Spec H_{\tau^\vee}^{(p)}(z^{-1})=\Spec H_\tau^{(p)}(z).
\]
从而 $R_p$ 在合法 $Q$-word 的二面体轨道上为常数，并与 lift 的选择无关。
\end{proposition}

\begin{proof}
令 $(Dx)_i=(-1)^ix_i$。位移一项在 $D$ 共轭下变号，位移二项不变，故
$DA_\tau D=-A_{-\tau}$。同时 $D$ 把 $\mathcal B_z^{(p)}$ 送到
$\mathcal B_{(-1)^pz}^{(p)}$。单位圆在相位映射下不变，因此平方谱边相同。

循环平移与反射分别由 $(S_kx)_i=x_{i+k}$ 和 $(\mathcal Rx)_i=x_{-i}$ 实现。直接代入得
\[
 A_{\operatorname{rot}_k\tau}S_k=S_kA_\tau,
 \qquad A_{\tau^\vee}\mathcal R=\mathcal RA_\tau.
\]
$S_k$ 保持 $\mathcal B_z^{(p)}$，$\mathcal R$ 把它送到
$\mathcal B_{z^{-1}}^{(p)}$，结论随即成立。
\end{proof}

若 $\tau$ 的最小周期为 $q$，但在 $p=kq$ 的胞元中展示，则由 $q$ 位平移进一步分解得到
\begin{equation}
 \Spec H_\tau^{(p)}(z)=\bigcup_{w^k=z}\Spec H_\tau^{(q)}(w).
 \label{eq:zone-folding}
\end{equation}
故胞元重复不改变 Bloch 谱边。下文把最小周期称为本原周期，把所选胞元长度称为展示周期。

\subsection{有限 Bloch 分解}
\label{sec:finite-bloch}

设 $\tau$ 的展示周期为 $p$ 且 $n=pL$。把坐标分成 $L$ 个胞元。胞元平移是酉算子，
其本征值恰为满足 $z^L=\alpha$ 的相位；对应本征空间内 $X_r=z^rX_0$。因为 Hamilton
规范算子与胞元平移可交换，每个本征空间都不变，其限制就是 $p\times p$ Hermitian
fiber $H_\tau^{(p)}(z)$。

\begin{proposition}
\label{prop:finite-bloch}
对任意 $L\geq1$ 与 $\alpha\in\{\pm1\}$，
\begin{equation}
 A_{pL,\alpha}(\tau)\simeq\bigoplus_{z^L=\alpha}H_\tau^{(p)}(z),
 \label{eq:finite-bloch-sum}
\end{equation}
并且
\begin{equation}
 \rho(A_{pL,\alpha}(\tau))=max_{z^L=\alpha}\rho(H_\tau^{(p)}(z)).
 \label{eq:finite-radius-max}
\end{equation}
\end{proposition}

\begin{proof}
胞元平移的本征空间两两正交，每个维数为 $p$，共有 $L$ 个，维数总和为 $pL$，故没有
遗漏任何本征值或重数。$|z|=1$ 时，所有跨胞元项与其共轭配对，所以 fiber 为 Hermitian。
谱半径公式由正交直和立即得到。
\end{proof}

\begin{figure}[t]
 \centering
 \begin{tikzpicture}[x=1.25cm,y=1cm,>=Latex,
 cell/.style={draw,rounded corners=2pt,minimum width=1.0cm,minimum height=.7cm,
              fill=blue!5,font=\small}]
 \node[cell] (c0) at (0,0) {$\tau$};
 \node[cell] (c1) at (1.35,0) {$\tau$};
 \node[cell] (c2) at (2.70,0) {$\tau$};
 \node (dots) at (4.0,0) {$\cdots$};
 \node[cell] (cl) at (5.25,0) {$\tau$};
 \draw[->] (c0)--node[above,font=\scriptsize] {$S$}(c1);
 \draw[->] (c1)--(c2);
 \draw[->] (c2)--(dots);
 \draw[->] (dots)--(cl);
 \draw[->,bend right=42,red!70!black]
   (cl.north) to node[above=5pt,font=\small] {seam $\alpha$} (c0.north);
 \draw[decorate,decoration={brace,mirror,amplitude=4pt}]
   (-.55,-.55)--node[below=5pt,font=\small] {$L$ 个胞元，$n=pL$}(5.8,-.55);
 \node[align=center,font=\small] at (2.65,-1.85)
 {$SX=zX$ 且 $S^L=\alpha I$\\[1pt]
  $\Longrightarrow\quad z^L=\alpha,\qquad
  A_{pL,\alpha}\simeq\displaystyle\bigoplus_{z^L=\alpha}H_\tau^{(p)}(z)$};
\end{tikzpicture}

 \caption{由 $L$ 个周期胞元构成的有限环。局部系数按胞重复，seam 携带 Hamilton
 holonomy $\alpha$，胞元平移的对角化给出允许相位 $z^L=\alpha$。}
 \label{fig:finite-bloch}
\end{figure}

有限直和把周期符号问题化成一族有限 Hermitian 矩阵。若要继续约化，就必须在每个 fiber
内部找到对称。系数 word 在半周期平移下反转，恰好能够迫使谱关于零对称。
